\documentclass[preprint,12pt]{elsarticle}
\usepackage{amssymb}
\usepackage{amsmath}
\usepackage{graphicx}
\usepackage[table]{xcolor}
\usepackage{booktabs}
\usepackage{siunitx}
\usepackage{subcaption}
\usepackage{hyperref}
\usepackage{geometry}
\graphicspath{{figures/}}
% vlastní makra
\newcommand{\R}{\mathbb{R}}
\newcommand{\vect}[1]{\boldsymbol{#1}}
\DeclareMathOperator{\argmin}{arg\,min}
\newtheorem{theorem}{Theorem}
\newtheorem{lemma}[theorem]{Lemma}
\journal{Sensors and Actuators A}

\begin{document}
\begin{frontmatter}
\title{Fibre Bragg Grating Sensors for Wrist Rehabilitation Monitoring\tnoteref{t1}}
\tnotetext[t1]{This work was supported by grant XY-123.}

\author[cvut]{Jan Nov\'ak\corref{cor1}}
\ead{jan.novak@fel.cvut.cz}
\author[cvut,vsb]{Petra Dvo\v{r}\'akov\'a}
\ead{petra.dvorakova@vsb.cz}
\author[vsb]{John Smith}
\cortext[cor1]{Corresponding author}

\affiliation[cvut]{organization={Faculty of Electrical Engineering, Czech Technical University in Prague},
            addressline={Technická 2},
            city={Prague},
            postcode={166 27},
            country={Czech Republic}}
\affiliation[vsb]{organization={VŠB -- Technical University of Ostrava},
            city={Ostrava},
            country={Czech Republic}}

\begin{abstract}
We present a wearable sensing system based on fibre Bragg gratings (FBG) for
monitoring wrist flexion during rehabilitation. The sensor achieves a resolution
of \SI{0.5}{\degree} and was validated on 12 participants.
\end{abstract}

\begin{highlights}
\item FBG sensor embedded in a flexible glove
\item Resolution of 0.5 degree
\end{highlights}

\begin{keyword}
fibre Bragg grating \sep rehabilitation \sep wearable sensors \sep optical fibre
\PACS 42.81.Pa
\end{keyword}
\end{frontmatter}

\section{Introduction}
\label{sec:intro}
Optical fibre sensors are immune to electromagnetic interference \citep{kersey1997,othonos1997}.
As shown by \citet{kersey1997}, FBGs can measure strain. % komentář v těle
Let $\vect{x}\in\R^n$ and consider $\argmin_{\vect{x}} f(\vect{x})$.

\section{Methods}
\subsection{Sensor design}
The sensor is shown in Fig.~\ref{fig:sensor}.
\begin{figure}[ht]
  \centering
  \includegraphics[width=0.8\textwidth]{sensor}
  \caption{Schematic of the FBG sensor.}
  \label{fig:sensor}
\end{figure}

\begin{theorem}\label{thm:a}
The Bragg wavelength shifts linearly with strain.
\end{theorem}
\begin{lemma}
Temperature compensation is needed.
\end{lemma}

\begin{table}[ht]
\centering
\caption{Measured parameters.}\label{tab:par}
\begin{tabular}{lcc}
\toprule
Parameter & Value & Unit\\
\midrule
Sensitivity & 1.2 & pm/$\mu\varepsilon$\\
Resolution & 0.5 & deg\\
\bottomrule
\end{tabular}
\end{table}

\begin{equation}
\Delta\lambda_B = \lambda_B (1-p_e)\varepsilon
\label{eq:bragg}
\end{equation}

\section{Results}
Results are summarized in Fig.~\ref{fig:res} and Table~\ref{tab:par}, see also \url{https://example.org}.
\begin{figure}[ht]
\centering
\begin{subfigure}{0.45\textwidth}
\includegraphics[width=\linewidth]{results}
\caption{A}
\end{subfigure}
\begin{subfigure}{0.45\textwidth}
\includegraphics[width=\linewidth]{figures/results.png}
\caption{B}
\end{subfigure}
\caption{Results.}\label{fig:res}
\end{figure}

\section{Conclusions}
FBG sensors are suitable for rehabilitation.

\section*{Acknowledgements}
We thank the volunteers who participated in the study.

\appendix
\section{Calibration data}
Additional data are given here.

\bibliographystyle{elsarticle-num}
\bibliography{refs}
\end{document}
